\documentclass{article}

\textwidth=6.5in
\textheight=8.5in
\voffset=-54pt
\oddsidemargin=0in
\evensidemargin=0in
\setlength{\parskip}{8pt}
\setlength{\parindent}{0pt}

\usepackage{workbook}

\usepackage[]{handout}

\begin{document}

\begin{center}
  \large\textsc{Problem Set 26 - Antiderivatives}
  
      \pspace{12pt}
  \begin{problemonly}\large Name: \rule{5cm}{0.15mm}\\ \end{problemonly}
\end{center}

\begin{grading}
\begin{enumerate}
    \item 10 points
    \item 10 points
    \item 10 points
    \item 8 points
  
\end{enumerate}
Total: 38 points.
\end{grading}

\begin{enumerate}
\item 

\begin{grading}
10 points. 2 per part. If they forget $+C$ on multiple parts, only deduct a half point once.
\end{grading}

  Find the following indefinite
  integrals.  Remember that an indefinite integral is always a family of
  functions (practically speaking, this means your answers should include
  ``$ + C$''). For help, read pp. 783-786 in Gottlieb, focusing on Example 25.1.

  \probonly{}
    \begin{enumerate}
    \item
      \begin{problem}[\vfill]
        $\displaystyle \int (3x^3 +2x +\pi)\,dx$
      \end{problem}

      \begin{solution}
        $\dfrac{3}{4}x^4+x^2+\pi x+C$
      \end{solution}

    \item 
      \begin{problem}[\vfill]
        $\displaystyle \int 3x^{-1} \,dx$
      \end{problem}

      \begin{solution}
        $3\ln |x| + C$
      \end{solution}

    \item
      \begin{problem}[\vfill]
        $\displaystyle \int \frac{e}{\pi x^2}\,dx$
      \end{problem}

      \begin{solution}
        $\displaystyle \int \frac{e}{\pi x^2}\,dx = \frac{e}{\pi} \int
        x^{-2}\,dx = \boxed{- \frac{e}{\pi} x^{-1} + C}$. 
      \end{solution}

    \item
      \begin{problem}[\nstretch{1}]
        $\displaystyle \int \left( 3\sin t -\frac{3}{1+t^2} \right)\,dt$
      \end{problem}

      \begin{solution}
        $-3\cos t -3\arctan t + C$
      \end{solution}

    \item
      \begin{problem}[\vfill]
        $\displaystyle \int \frac{2\cos w}{3}\, dw$
      \end{problem}

      \begin{solution}
        $ \dfrac{2\sin w}{3} + C$ 
      \end{solution}

    \end{enumerate}    
  \probonly{}
  $\checkmark$ Remember that you can check your answers by
    differentiating them!

\item 

\begin{grading}
10 points. 2 for each subpart of (a), and 2 for part (b). They don't need $+C$ for part (a).
\end{grading}
  \begin{enumerate}
  \item
Find antiderivatives for the given functions. In other words, for each
    function $f$, find a function $F$ such that $F' = f$.

    \probonly{}
      \begin{enumerate}
      \item
        \begin{problem}[\vfill]
          $f(x) = e^{3x}\vphantom{\dfrac{1}{1}}$
        \end{problem}

        \begin{solution} % Jason Bland
          We can take $F(x) = \boxed{\frac{1}{3} e^{3x}}$ (plus any
          constant).
        \end{solution}

      \item
        \begin{problem}[\vfill]
          $f(x) = \dfrac{3}{e^x}$
        \end{problem}

        \begin{solution} % Jason Bland
          We can rewrite $f(x) = 3 e^{-x}$, so $F(x) = \boxed{-3 e^{-x}}$
          works.
        \end{solution}

      \item
        \begin{problem}[\vfill\newpage]
          $f(x) = \dfrac{1}{2x}$
        \end{problem}

        \begin{solution}
          We can rewrite $\displaystyle f(x) = \frac{1}{2} \cdot
          \frac{1}{x}$, so $F(x) = \boxed{\dfrac{1}{2} \ln |x|}$ is one
          antiderivative.
        \end{solution}

      \item
        \begin{problem}[\vfill]
          $f(x) = \dfrac{(x + 1)^2}{\sqrt{x}}$ \\
          \emph{Hint: expand the numerator so you can break up the 
          fraction into the sum of simpler pieces.}
        \end{problem}

        \begin{solution}
          We don't have a Quotient Rule for antiderivatives, so let's
          first expand this out: $\displaystyle f(x) = \frac{x^2 + 2x +
            1}{\sqrt{x}} = x^{3/2} + 2 x^{1/2} + x^{-1/2}$.  Then, one
          antiderivative of $f$ is $F(x) = \boxed{\frac{2}{5} x^{5/2} +
            \frac{4}{3} x^{3/2} + 2 x^{1/2}}$.
        \end{solution}

      \end{enumerate}
      \probonly{}

  \item
    \begin{problem}[\nstretch{2}]
      Check your answers to {{{{\#2}}(a)}} by
      differentiating them and write your work below.
    \end{problem}

    \begin{solution}
      \begin{enumerate}
      \item The derivative of $F(x) = \frac{1}{3} e^{3x}$ is $\frac{1}{3}
        \cdot 3 e^{3x} = e^{3x}$.
      \item The derivative of $F(x) = - 3 e^{-x}$ is $- 3 \left( - e^{-x}
        \right) = 3 e^{-x}$.
      \item The derivative of $F(x) = \frac{1}{2} \ln |x|$ is $\frac{1}{2}
        \cdot \frac{1}{x} = \frac{1}{2x}$.
      \item The derivative of $F(x) = \frac{2}{5} x^{5/2} +
            \frac{4}{3} x^{3/2} + 2 x^{1/2}$ is $x^{3/2} + 2 x^{1/2} +
            x^{-1/2}$, which can be rewritten as $\frac{x^2 + 2x +
              1}{\sqrt{x}}$. 
      \end{enumerate}
    \end{solution}

  \end{enumerate}

\item 
\begin{grading}
10 points, as follows: 4/4/2.
\end{grading}

  \begin{enumerate}

  \item
    \begin{problem}[\vfill]
      Suppose the velocity of an object at time $t$ is $v(t) =
      \frac{1}{1+t^2}$ miles per hour, where $t$ is measured in hours.
      What is the net change in the object's position from $t=0$ to $t=1$?
      What is the total distance the object travels between $t = 0$ and $t
      = 1$?
    \end{problem}

    \begin{solution}
      The object's net change in position from $t = 0$ to $t = 1$ is
      $\displaystyle \int_0^1 \frac{1}{1 + t^2}\,dt = \arctan t \Big|_0^1
      = \boxed{\frac{\pi}{4} \text{ miles}}$.

      Since the rate of change of distance traveled is speed and the
      object's speed at time $t$ is $|v(t)|$, the distance traveled is
      $\displaystyle \int_0^1 |v(t)|\,dt$.  In this case, the object's
      velocity $\frac{1}{1 + t^2}$ is always positive, so its speed is the
      same as its velocity; therefore, the total distance traveled is also
      $\boxed{\frac{\pi}{4} \text{ miles}}$. 
    \end{solution}

  \item
    \begin{problem}[\vfill]
      Suppose the velocity of another object at time $t$ is $v(t) =
      5t(t-1)$ meters per second, where $t$ is measured in seconds.
      What is the net change in the object's position from $t=0$ to $t=2$?
      What is the total distance the object travels between $t = 0$ and $t
      = 2$?
    \end{problem}

    \begin{solution}
      The object's net change in position from $t = 0$ to $t = 2$ is
      \begin{align*}
        \int_0^2 5t (t - 1)\,dt &= 5 \int_0^2 (t^2 - t)\,dt \\
        &= 5 \left( \left. \frac{t^3}{3} - \frac{t^2}{2} \right|_0^2
        \right) \\
        &= \boxed{\frac{10}{3} \text{ meters}}
      \end{align*}
      This object's velocity is negative on $(0, 1)$ and positive on $(1,
      2)$\footnote{You can see this by making a sign chart for $v(t)$, or
        by graphing it: you can graph it by hand since it's a quadratic.},
      so the total distance it travels between $t = 0$ and $t = 2$ is
      different from its net change in position.  To get the total
      distance traveled, we need to integrate the object's speed, which is
      $|v(t)|$.  On $(0, 1)$, $v(t)$ is negative, so $|v(t)| = -v(t) =
      -5t(t-1)$; on $(1, 2)$, $v(t)$ is positive, so $|v(t)| = v(t) =
      5t(t-1)$.  So, the total distance the object travels between $t = 0$
      and $t = 2$ is:
      \begin{align*}
        \int_0^1 -5t(t-1)\,dt + \int_1^2 5t(t-1)\,dt &= -5 \int_0^1 (t^2 -
        t)\,dt + 5 \int_1^2 (t^2 - t)\,dt \\
        &= -5 \left( \left. \frac{t^3}{3} - \frac{t^2}{2} \right|_0^1
        \right) + 5 \left( \left. \frac{t^3}{3} - \frac{t^2}{2} \right|_1^2
        \right) \\
        &= -5 \left( - \frac{1}{6} \right) + 5 \left( \frac{5}{6} \right)
        \\
        &= \boxed{5 \text{ meters}}
      \end{align*}
      (You may also want to look back at
        {{\#5}} on the worksheet for another example of the
        distinction between ``change in position'' and ``distance
        traveled''.)
    \end{solution}

  \item
    \begin{problem}[\vfill]
      When is net change in position equal to total distance traveled?
      (That is, what must be true about the velocity function?)
    \end{problem}

    \begin{solution}
      Net change in position is equal to total distance traveled if the
      velocity is always positive.  (This should make sense intuitively:
      if the object is always traveling in the positive direction, then
      its net change in position will simply be the distance it traveled.)

      \emph{Note:} If the object is always traveling in the negative
      direction, then its net change in position will simply be the
      negative of the distance traveled. 
    \end{solution}
  \end{enumerate}

\iffalse
\item 

An economist has been tracking housing prices in Cambridge for the past
  several years.  He estimates that the \emph{rate} at which the median
  price of a 1-bedroom house has been increasing over the past several
  years is an exponential function.  At time $t = 0$, corresponding to
  January 2007, the median price was increasing at a rate of \$15,000 per
  year.  At time $t = 1$, corresponding to January 2008, the median price
  was increasing at a rate of \$15,750 per year.
  \begin{enumerate}
  \item

    \begin{problem}
      Write an exponential function $r(t)$ which models the \emph{rate} at
      which the median price is increasing over time, in dollars per
      year.  
    \end{problem}

    \begin{solution}
      We are looking for an exponential function $r(t)$ with $r(0) =
      15,000$ and $r(1) = 15,750$.  Remember that exponential functions
      have a constant \textbf{percent} change per year.\footnote{They also
        have a constant percent change per month, or per day, or per any
        other unit of time.}

      In the first year, the percent change in $r(t)$ was $\frac{15750 -
        15000}{15000} = 5\%$.  So, our exponential function $r(t)$ must
      increase by 5\% per year: $\boxed{r(t) = 15,000 (1.05)^t}$.
    \end{solution}

  \calcitem
    \begin{problem}
      At time $t = 0$, the median price of a 1-bedroom house in Cambridge
      was \$300,000.  Assuming the economist's trend held, find the median
      price of a 1-bedroom house in Cambridge at $t = 2$.  Please give an
      exact answer, and then use a calculator to give a decimal
      approximation. 
    \end{problem}

    \begin{solution}
      The house's price at $t = 2$ is
      \begin{align*}
        300,000 + \int_0^2 r(t)\,dt &= 300,000 + 15,000 \int_0^2
        1.05^t\,dt \\
        &= 300,000 + 15,000 \left( \left. \frac{1.05^t}{\ln 1.05}
          \right|_0^2 \right) \\
        &= \boxed{300,000 + 15,000 \left( \frac{1.05^2 - 1 }{\ln 1.05}
          \right)} \\
        & \boxed{\approx 331512.50}
      \end{align*}
    \end{solution}

  \end{enumerate}  
\fi

\pspace{\newpage}

\item 

\begin{grading}
8 points. 3 for part (a), 5 for part (b).
\end{grading}

 \textbf{Reflect Back.}

  \begin{enumerate}
  \item Decide whether each of the following statements is true or false. Write 1-2 sentences explaining your choice.

  \probonly{\raggedcolumns}
    \begin{enumerate}
    \item
      \begin{problem}
        $\displaystyle \int \cos u\,du = \sin u + C$
      \end{problem}

      \begin{solution} % Jason Bland
        True, because the derivative of $\sin u$ is $\cos u$.
      \end{solution}

    \item
      \begin{problem}
        $\displaystyle \int \cos(x + 7)\,dx = \sin(x + 7) + C$
      \end{problem}

      \begin{solution} % Jason Bland
        True, because the derivative of $\sin(x + 7)$ is $\cos(x + 7)
        \cdot 1$. 
      \end{solution}

    \item
      \begin{problem}
        $\displaystyle \int \cos (3x + 7)\,dx = \sin(3x + 7) + C$
      \end{problem}

      \begin{solution} % Jason Bland
        False, because the derivative of $\sin(3x + 7)$ is $\cos(3x + 7)
        \cdot 3$.
      \end{solution}

    \item
      \begin{problem}
        $\displaystyle \int \cos (3x + 7)\,dx = \frac{1}{3} \sin (3x + 7)
        + C$
      \end{problem}

      \begin{solution} % Jason Bland
        True, because the derivative of $\frac{1}{3} \sin(3x + 7)$ is
        $\frac{1}{3} \cos(3x + 7) \cdot 3$.
      \end{solution}

    \item
      \begin{problem}
        $\displaystyle \int \cos (x^2)\,dx = \sin (x^2) + C$
      \end{problem}

      \begin{solution} % Jason Bland
        False, because the derivative of $\sin(x^2)$ is $\cos(x^2) \cdot
        2x$. 
      \end{solution}

    \item
      \begin{problem}
        $\displaystyle \int \cos (x^2)\,dx = \frac{1}{2x} \sin (x^2) + C$
      \end{problem}

      \begin{solution} % Jason Bland
        False, because we need to use either the Product or Quotient
        Rule to evaluate the derivative of $\frac{1}{2x} \sin (x^2) =
        \frac{1}{2} x^{-1} \sin(x^2)$; the derivative is $\frac{1}{2}
        \left[ - x^{-2} \sin(x^2) - x^{-2} \cos(x^2) \cdot 2x \right]$.
      \end{solution}

    \end{enumerate}
    \probonly{}

    \probonly{\emph{Hint:} Exactly three are true.}

    \pspace{\newpage}

  \item

    \note{When we do integration by substitution, it will be really
      important for students to know what integrals they know.}

    \begin{problem}[\vfill]
      Five of the indefinite integrals below are ones you should 
      know.  (The others are ones that you are not expected to be able to
      find yet.)  Pick out the five that you know, and find them.
      \[ \int \sin x\,dx \quad \int \tan x\,dx \quad \int e^x\,dx \quad \int
      \ln x\,dx \quad \int \frac{1}{x}\,dx \]
      \[ \int \frac{1}{1 + x^2}\,dx \quad \int \frac{1}{1 - x^2}\,dx \quad
      \int \frac{1}{\sqrt{1 + x^2}}\,dx \quad \int \frac{1}{\sqrt{1 -
          x^2}}\,dx \]
    \end{problem}

    \begin{solution}
      The five we know are:
      \begin{itemize}
      \item $\displaystyle \int \sin x\,dx = - \cos x + C$
      \item $\displaystyle \int e^x \,dx = e^x + C$
      \item $\displaystyle \int \frac{1}{x}\,dx = \ln |x| + C$
      \item $\displaystyle \int \frac{1}{1 + x^2}\,dx = \arctan x + C$
      \item $\displaystyle \int \frac{1}{\sqrt{1 - x^2}}\,dx = \arcsin x +
        C$ (or $-\arccos x + C$)
      \end{itemize}
    \end{solution}

  \end{enumerate}


\end{enumerate}

\spaceonly{\psreflection{}}

\end{document}